\chapter*{Contribuciones Personales}
\label{cap:contribucionesPersonales}
\addcontentsline{toc}{chapter}{Contribuciones Personales}

\section*{Miguel Xuclá Herrero}
Durante las fases iniciales del proyecto, Miguel contribuyó a la comprensión y manejo del hardware empleado, investigando el uso correcto del casco EEG y el procedimiento adecuado para la reubicación y configuración de sus electrodos. Asimismo, llevó a cabo un análisis exhaustivo de diversos \textit{datasets} públicos orientados a interfaces cerebro-computador, estudiando a fondo su estructura y prestando especial atención a la duración temporal de los \textit{trials} y a los periodos de descanso entre estos (véase \hyperref[subsec:estructuraGrabacion]{estructura grabación}). En paralelo a estas tareas, prestó apoyo en la etapa de investigación de inteligencia artificial, encargándose de evaluar y comparar diferentes modelos de clasificación para determinar su viabilidad.

En etapas posteriores, durante el desarrollo del \textit{pipeline} en tiempo real, Miguel desempeñó un rol fundamental en el tratamiento de las predicciones (véase el Capítulo \ref{cap:aplicacionAVideojuegos}). Investigó, ideó y probó empíricamente diversos filtros de decisión, realizando un análisis del flujo de predicciones para evaluar la estabilidad y eficacia del sistema antes y después de aplicar dichas técnicas de suavizado.

Fue el encargado de pensar, diseñar y revisar a conciencia los videojuegos experimentales utilizados en las pruebas. Vinculado a este desarrollo, Miguel ideó e implementó algorítmicamente los distintos métodos que conforman los filtros de ayuda integrados en los propios juegos, desarrollando la lógica matemática necesaria para asistir al usuario y compensar el ruido inherente a la señal.

En las últimas fases del proyecto, colaboró estrechamente en el diseño del protocolo definitivo seguido durante los experimentos en tiempo real (véase el Capítulo \ref{cap:experimentoTiempoReal}), participando también en la creación del formulario de evaluación que respondieron los sujetos. Finalmente, participó activamente en el análisis de las grabaciones en vídeo generadas durante las sesiones de juego, encargándose de la recolección de métricas y del posterior análisis e interpretación de los resultados obtenidos.

\section*{Jorge Sáez Salto}
Jorge llevó a cabo una investigación inicial sobre técnicas de extracción de características de señales (ICA, Fourier, Wavelet), estudio de los distintos paradigmas EEG, rangos de frecuencia del cerebro humano, filtros de frecuencia, búsqueda de repositorios con \textit{datasets}, modelos de clasificación especializados en este campo, y posteriormente técnicas de detección y corrección de \textit{trials} ruidosos y métodos de filtrado de predicciones de modelos de clasificación.

En las primeras fases del proyecto, Jorge se dedicó especialmente a la búsqueda y preparación en Python de un modelo de clasificación, siendo MiRepNet la opción finalmente escogida. Asimismo, se encargó de la resolución de dependencias de librerías y entornos, logrando realizar el primer \textit{test} de funcionamiento del modelo.

Jorge desempeñó un rol multifunción durante las fases siguientes: participó en el desarrollo de la aplicación, concretamente en los módulos de configuración de electrodos y experimento visual con \textit{pygame}, y el \textit{SDK} del \textit{BrainAccess MIDI} (véase el Capítulo~\ref{cap:aplicacionBCI}), y asistió a dos de las cuatro sesiones de grabación piloto, supervisando el transcurso de las mismas. Posteriormente, se encargó de la integración de las grabaciones obtenidas con el modelo de clasificación, llevando a cabo numerosas pruebas con distintos filtros y normalizaciones hasta su correcta implementación, así como del desarrollo de la fase de entrenamiento para la extracción de la matriz de alineamiento euclidiano y el modelo preentrenado previo al experimento en tiempo real.

Realizó investigaciones y pruebas sobre detección de \textit{trials} ruidosos e interpolación de canales, estudiando métodos como FASTER y Autoreject, y participó en la toma de decisiones sobre los criterios de detección de \textit{trials} ruidosos a aplicar en la arquitectura final.

En las últimas fases del proyecto, Jorge tuvo un papel importante en el desarrollo del experimento en tiempo real, realizando simulaciones \textit{offline}, prueba de procesos intérpretes, e implementación y pruebas con los distintos filtros de decisión, y colaborando en la creación de los sistemas de ayuda para los diferentes juegos, así como su integración en la aplicación del sistema BCI. Por último, se encargó de redactar el protocolo, preparar el entorno y realizar las grabaciones en vídeo de los tres experimentos con videojuegos, así como participar en la extracción de datos y métricas de los mismos y la elaboración de conclusiones (véase el Capítulo~\ref{cap:experimentoTiempoReal}).

\section*{Julián Reguera Peñalosa}
Julián se encargó, en una primera fase del proyecto, de realizar una revisión bibliográfica orientada a comprender el estado actual de las interfaces cerebro-computador y las posibilidades reales de aplicación de este tipo de sistemas. Para ello, llevó a cabo una investigación sobre los principales paradigmas BCI, las limitaciones habituales presentes en este tipo de tecnologías y los métodos más utilizados para mejorar el rendimiento de los sistemas basados en señales EEG. Asimismo, estudió diferentes librerías y herramientas empleadas en el ámbito del procesamiento de señal y aprendizaje automático aplicado a BCI, así como conceptos fisiológicos y técnicos relacionados con la adquisición y tratamiento de señales electroencefalográficas. Durante esta etapa también analizó los rendimientos esperados en tareas de imaginación motora y los factores que condicionan la precisión de estos sistemas. Tras comparar distintas alternativas, participó en la selección de la librería MNE como base principal para el procesamiento EEG, profundizando posteriormente en su funcionamiento interno y en las herramientas que proporciona para el tratamiento de señales biomédicas.

Posteriormente, asumió el estudio de la librería asociada al dispositivo EEG \textit{BrainAccess MIDI}, así como de toda la documentación y API proporcionada por el fabricante. Debido a la escasa información disponible en la documentación oficial, fue necesario realizar un análisis detallado del código fuente de la propia librería, especialmente de la clase encargada de la adquisición EEG y de la gestión de buffers internos, con el objetivo de comprender el funcionamiento real de las funciones de configuración, captura y lectura de datos. Como resultado de esta investigación, adquirió un conocimiento profundo sobre el proceso de adquisición de señales EEG desde Python y sobre las particularidades técnicas del dispositivo utilizado en el proyecto.

Durante las primeras fases de desarrollo colaboró junto a Jorge Sáez Salto en la preparación del entorno de trabajo común del proyecto. Para ello desarrolló distintos scripts de instalación y ejecución orientados a simplificar la configuración del entorno software, automatizar dependencias y facilitar el despliegue de la aplicación al resto de integrantes del equipo.

Más adelante llevó a cabo una investigación específica sobre métodos de interpolación espacial aplicados a señales EEG, estudiando especialmente los algoritmos empleados de forma habitual en herramientas de referencia como MNE. A partir de este trabajo desarrolló la primera versión funcional de la aplicación EEG del proyecto, la cual permitía configurar el dispositivo, comprobar el estado e impedancia de los electrodos y realizar grabaciones básicas de señal en tiempo real.

Tras la selección del modelo MiRepNet como arquitectura principal para el sistema de clasificación, participó junto al resto del grupo en un análisis exhaustivo tanto de la estructura interna del modelo como de su funcionamiento completo. Debido a la limitada documentación disponible en el repositorio oficial, fue necesario estudiar detalladamente el código fuente del proyecto y analizar el comportamiento real de las distintas etapas del pipeline de procesamiento. Este análisis permitió identificar discrepancias relevantes entre la implementación y la descripción presentada en el artículo científico original, incluyendo erratas relacionadas con el orden de las etapas de preprocesado y con la utilización de determinados filtros que, aunque no aparecían correctamente reflejados en la publicación, resultaban necesarios para reproducir adecuadamente el funcionamiento del modelo.

Dentro de este proceso también llevó a cabo un estudio detallado de las técnicas de interpolación espacial y del método de alineamiento euclídeo empleados por MiRepNet, con el objetivo de comprender su fundamento matemático y su impacto sobre la representación de las señales EEG. A partir de este análisis investigó la viabilidad de adaptar dichas técnicas a un entorno de inferencia en tiempo real, evaluando diferentes estrategias de implementación y las limitaciones prácticas asociadas a este tipo de procesamiento. En particular, estudió el comportamiento de las matrices de transformación obtenidas mediante alineamiento euclídeo y las diferencias existentes entre sujetos, entre distintas sesiones de grabación de un mismo sujeto y entre los conjuntos de entrenamiento y validación, analizando cómo estas variaciones podían afectar al rendimiento del sistema.

Además, participó en la definición del protocolo experimental empleado durante las fases preliminares de adquisición de señal, colaborando en la planificación de las sesiones, la validación de los paradigmas utilizados y la evaluación de la calidad de las grabaciones obtenidas. También participó activamente en las sesiones de captura del experimento piloto y fue el principal responsable de la realización de las grabaciones definitivas utilizadas posteriormente tanto en la evaluación de técnicas de preprocesado como en el entrenamiento y validación del modelo de clasificación.

En relación con el procesamiento de señal, colaboró en la investigación y evaluación de diferentes técnicas de filtrado, normalización y preparación de señales EEG para tareas de imaginación motora, estudiando el impacto de cada una de ellas sobre el rendimiento del sistema final. Asimismo, participó en el estudio de la posible adaptación del modelo MiRepNet al problema concreto abordado en el proyecto y en la evaluación de distintas configuraciones experimentales.

Finalmente, dentro del sistema de tiempo real desarrollado en la aplicación EEG, diseñó e implementó la arquitectura software basada en procesos independientes utilizada por el sistema, así como los mecanismos de comunicación y sincronización entre procesos. También desarrolló una parte importante de la cadena de adquisición y procesamiento en tiempo real, incluyendo la captura continua de datos EEG, la implementación de buffers circulares para el manejo eficiente de ventanas temporales, la aplicación de filtros y normalizaciones en línea y la integración del modelo de clasificación dentro del flujo completo de inferencia en tiempo real.
